% Generated draft: requires mathematical and source review.
\documentclass[11pt]{article}
\usepackage[margin=25mm]{geometry}
\usepackage{fontspec}
\setmainfont{Latin Modern Roman}
\usepackage{amsmath,amssymb,mathtools,unicode-math}
\setmathfont{Latin Modern Math}
\usepackage{enumitem,needspace,xurl,hyperref}
\setlength{\parindent}{0pt}
\setlength{\parskip}{5pt}
\setlength{\emergencystretch}{3em}
\newcommand{\N}{\mathbb{N}}
\newcommand{\Z}{\mathbb{Z}}
\newcommand{\Q}{\mathbb{Q}}
\newcommand{\R}{\mathbb{R}}
\newcommand{\C}{\mathbb{C}}
\newcommand{\F}{\mathbb{F}}
\newcommand{\A}{\mathbb{A}}
\newcommand{\PP}{\mathbb{P}}
\newcommand{\E}{\mathbb{E}}
\newcommand{\eps}{\varepsilon}
\newcommand{\dd}{\,\mathrm{d}}
\begin{document}
\section*{UnsolvedMath ID: 30003364}
\textbf{Source problem code:} catalogue entry
\subsection{Definitions and mathematical statement}
For a positive integer \(n\), let \(\mathbb{CP}^{n}\) denote complex projective \(n\)-space, regarded as a smooth compact manifold without boundary of real dimension \(2n\). Under the interpretation that ``positively curved'' means positive sectional curvature, define
\[
\mathcal R_{\sec>0}(\mathbb{CP}^{n})
=\left\{g\in\Gamma^{\infty}(\operatorname{Sym}^{2}T^{*}\mathbb{CP}^{n}):\begin{array}{l}g\text{ is positive definite, and}\\ \sec_g(\sigma)>0\text{ for every }x\in\mathbb{CP}^{n}\\ \text{and every real two-dimensional subspace }\sigma\subseteq T_x\mathbb{CP}^{n}\end{array}\right\}.
\]
One conventional moduli space is
\[
\mathcal M_{\sec>0}(\mathbb{CP}^{n})
=\mathcal R_{\sec>0}(\mathbb{CP}^{n})/\operatorname{Diff}(\mathbb{CP}^{n}),
\]
where the full smooth diffeomorphism group acts by pullback, \(\phi\cdot g=\phi^{*}g\). For this convention, the metric space carries the \(C^{\infty}\) topology and the moduli space carries the quotient topology; no additional identification by constant rescaling is imposed.

\textbf{Source question, with unresolved conventions.} What can be said about the moduli space of positively curved metrics when the sphere \(S^{d}\) is replaced by complex projective space? With the conventions just displayed, this asks for information about \(\mathcal M_{\sec>0}(\mathbb{CP}^{n})\).

The displayed definitions provide an explicit conventional interpretation, not conventions established by the supplied record. The record does not specify the curvature condition, the equivalence relation, the topology, a normalization, the range of projective dimensions, or a particular invariant or theorem to determine. Thus \(n\geq1\) indexes the spaces defined here and is not a recovered quantifier from the source. No relation between \(d\) and \(n\) is supplied. The original wording also includes the sphere case, whereas the cleaned statement focuses on its replacement by complex projective space. This is an open-ended research question, not a supplied conjectural equality, classification, or quantified assertion.
\subsection{Short English statement}
The question asks what can be learned about moduli spaces of positively curved Riemannian metrics on complex projective spaces, in comparison with the sphere setting. Positive sectional curvature and quotienting by all diffeomorphisms give one conventional formulation, but the supplied record leaves these conventions, the dimensions, and the precise desired conclusion unresolved.
\subsection{Sources}
\begin{itemize}
\item Local UnsolvedMath registry: permanent ID \(30003364\), title \emph{Moduli of Positively Curved Metrics on Projective Spaces}, source problem code \texttt{catalogue entry}, collection \emph{Oberwolfach Reports Open Problems}. No individual proposer is supplied; the proposed year is recorded as 2017.
\item Supplied source citation: \emph{Mini-Workshop: Spaces and Moduli Spaces of Riemannian Metrics (2018)}, Oberwolfach Reports, pp.\ 133--166, DOI \texttt{10.4171/owr/2017/3}. \url{https://doi.org/10.4171/owr/2017/3}. The supplied literature-source label is ``Oberwolfach discussion of positively curved moduli.'' The citation's parenthetical year 2018 differs from the recorded proposed year and the year appearing in the DOI; this discrepancy is unresolved.
\item Supplied literature source: ``Kreck--Stolz nonconnected positive-curvature moduli.'' \url{https://doi.org/10.2307/2152742}.
\item Supplied literature source: ``Moduli spaces of positive curvature metrics in dimension four and beyond.'' \url{https://doi.org/10.1007/s00209-026-03947-3}.
\item The supplied literature assessment was checked on \texttt{2026-08-21}. It labels the problem open and reports special-case results on components and higher homotopy, including recent work on projective spaces, while reporting no general classification or homotopy description. This is supplied dated curation derived from the OWR record, not an independently verified assessment. The background explicitly calls for independent verification of time-sensitive status and cited claims.
\item The registry labels the statement \texttt{corrected\_verified}, explaining that notation, typography, and sentence structure were repaired without changing mathematical content. This supplied editorial label does not establish the missing mathematical conventions or constitute independent mathematical verification.
\end{itemize}
\end{document}
